% pdf-name: Platform Engineer - DevOps
% Variant "platform": Platform Engineer / DevOps with a developer background, full remote.
% \afaire{...} marks a placeholder (printed in red): fill it or remove it before sending.

\newcommand{\cvtitle}{Platform Engineer / DevOps -- développeur Java -- 19 ans d'expérience}

\newcommand{\cvprofile}{Développeur Java devenu DevOps en automatisant les déploiements de mon propre projet (2017), puis responsable de la plateforme d'une application de 3 feature teams : CI/CD, Kubernetes, infrastructure as code, observabilité. Je partage mon temps à 50/50 entre la plateforme et le développement d'outils internes, avec une connaissance concrète des besoins des équipes de développement. Habitué au travail à distance avec des équipes en Inde et en Pologne. Recherche un poste de Platform Engineer en full remote.}

\newcommand{\cvskills}{
        CI/CD & GitHub Actions, Argo CD (GitOps), Jenkins, Helm \\
        \hline
        Infra & Kubernetes, Docker, Terraform, Ansible, Shell, \afaire{Cloud : AWS / GCP / Azure si applicable} \\
        \hline
        Monitoring & OpenTelemetry, Stack Elastic (ELK, APM), Grafana \\
        \hline
        Langages & Java, Kotlin, Python, TypeScript, JavaScript, HCL \\
        \hline
        Dev & Spring Boot, microservices, APIs REST, Angular, JUnit, Cucumber, Gatling, Playwright \\
        \hline
        Data & Oracle, PostgreSQL, Cassandra, Elasticsearch \\
        \hline
        Méthodes & Scrum, Kanban, BDD, TDD \\
        \hline
}

\newcommand{\cvcurrentrole}{DevOps / Technical Leader}

\newcommand{\cvcurrentintro}{
  {Conférences suivies : KubeCon, DevoxxFR.\\}
}

\newcommand{\cvkeypoints}{
    \begin{itemize}
      \item Plateforme et DevOps
        \begin{itemize}
          \item Déploiement continu de \afaire{N} microservices sur Kubernetes (GitHub Actions, ArgoCD, Helm) : \afaire{fréquence de mise en production avant / après, ex. mensuelle vers quotidienne}
          \item Création des environnements en infrastructure as code (Terraform, Ansible) : \afaire{délai de création d'un environnement avant / après}
          \item Observabilité de l'application (Stack Elastic, APM, OpenTelemetry) : \afaire{ex. alerting, tableaux de bord, réduction du temps de diagnostic}
          \item Suivi de la production et gestion des incidents : \afaire{périmètre, astreinte éventuelle}
        \end{itemize}
      \item Outillage développeur et encadrement technique
        \begin{itemize}
          \item Développement fullstack du "Release Tool" (Java/Angular), outil interne de suivi des mises en production utilisé par \afaire{N équipes}
          \item Encadrement technique de 3 feature teams et d'une équipe DevOps (Inde, Pologne, France)
          \item Définition et suivi de la roadmap technique, normes de développement du groupe
          \item Animation de "Leagues" et "Chapters" sur le software craftsmanship
        \end{itemize}
    \end{itemize}
}
